\documentclass[10pt,a4paper]{amsart}
\usepackage{amsmath,amssymb,amsthm,graphicx}
% Retained preamble for the Stacks Project chapter divisors.tex (Divisors).
% The chapter file pinned-original-source.tex carries only \input{preamble} on its first line, so no
% preamble text lives inside the chapter itself.  The Stacks Project uses one single shared preamble for
% every chapter; the snapshot stages it verbatim in sources/stacks-chow/preamble.tex.  The block below is
% copied mechanically, line for line, from that file's lines 14-18 (the Xy-pic setup) and lines 154-217
% (the native theorem declarations together with the mathematical macros).  Nothing was added, removed or
% redefined.  The wrapper class and the amsmath/amssymb/amsthm/graphicx base packages are supplied by the
% assembler, so the class/\externaldocument part of the shared preamble is not repeated here.
% For commutative diagrams we use Xy-pic
\usepackage[all]{xy}

% We use 2cell for 2-commutative diagrams.
\xyoption{2cell}
% Theorem environments.
%
\theoremstyle{plain}
\newtheorem{theorem}[subsection]{Theorem}
\newtheorem{proposition}[subsection]{Proposition}
\newtheorem{lemma}[subsection]{Lemma}

\theoremstyle{definition}
\newtheorem{definition}[subsection]{Definition}
\newtheorem{example}[subsection]{Example}
\newtheorem{exercise}[subsection]{Exercise}
\newtheorem{situation}[subsection]{Situation}

\theoremstyle{remark}
\newtheorem{remark}[subsection]{Remark}
\newtheorem{remarks}[subsection]{Remarks}

\numberwithin{equation}{subsection}

% Macros
%
\def\lim{\mathop{\mathrm{lim}}\nolimits}
\def\colim{\mathop{\mathrm{colim}}\nolimits}
\def\Spec{\mathop{\mathrm{Spec}}}
\def\Hom{\mathop{\mathrm{Hom}}\nolimits}
\def\Ext{\mathop{\mathrm{Ext}}\nolimits}
\def\SheafHom{\mathop{\mathcal{H}\!\mathit{om}}\nolimits}
\def\SheafExt{\mathop{\mathcal{E}\!\mathit{xt}}\nolimits}
\def\Sch{\mathit{Sch}}
\def\Mor{\mathop{\mathrm{Mor}}\nolimits}
\def\Ob{\mathop{\mathrm{Ob}}\nolimits}
\def\Sh{\mathop{\mathit{Sh}}\nolimits}
\def\NL{\mathop{N\!L}\nolimits}
\def\CH{\mathop{\mathrm{CH}}\nolimits}
\def\proetale{{pro\text{-}\acute{e}tale}}
\def\etale{{\acute{e}tale}}
\def\QCoh{\mathit{QCoh}}
\def\Ker{\mathop{\mathrm{Ker}}}
\def\Im{\mathop{\mathrm{Im}}}
\def\Coker{\mathop{\mathrm{Coker}}}
\def\Coim{\mathop{\mathrm{Coim}}}

% Boxtimes
%
\DeclareMathSymbol{\boxtimes}{\mathbin}{AMSa}{"02}

%
% Macros for moduli stacks/spaces
%
\def\QCohstack{\mathcal{QC}\!\mathit{oh}}
\def\Cohstack{\mathcal{C}\!\mathit{oh}}
\def\Spacesstack{\mathcal{S}\!\mathit{paces}}
\def\Quotfunctor{\mathrm{Quot}}
\def\Hilbfunctor{\mathrm{Hilb}}
\def\Curvesstack{\mathcal{C}\!\mathit{urves}}
\def\Polarizedstack{\mathcal{P}\!\mathit{olarized}}
\def\Complexesstack{\mathcal{C}\!\mathit{omplexes}}
% \Pic is the operator that assigns to X its picard group, usage \Pic(X)
% \Picardstack_{X/B} denotes the Picard stack of X over B
% \Picardfunctor_{X/B} denotes the Picard functor of X over B
\def\Pic{\mathop{\mathrm{Pic}}\nolimits}
\def\Picardstack{\mathcal{P}\!\mathit{ic}}
\def\Picardfunctor{\mathrm{Pic}}
\def\Deformationcategory{\mathcal{D}\!\mathit{ef}}

% --- appended 2026-10-12 (Chief): the project own definitions of the reference/slogan/history
% display environments, copied mechanically from the shared Stacks preamble lines 8-12
% (\usepackage{verbatim} plus the three \newenvironment lines, which the chapter preamble keeps
% verbatim from the project preamble). Nothing added, removed or redefined: these environments
% carry no mathematical content and the project itself typesets them as comments. Without these
% three lines a counted statement whose own text uses them fails with Environment slogan undefined.
\usepackage{verbatim}
\newenvironment{reference}{\comment}{\endcomment}
\newenvironment{slogan}{\comment}{\endcomment}
\newenvironment{history}{\comment}{\endcomment}
\usepackage{cleveref}
\crefname{originallemma}{original lemma}{original lemmas}
\Crefname{originallemma}{Original lemma}{Original lemmas}
\makeatletter
\newlabel{algebra-lemma-weakly-ass-zero-divisors}{{05C3}{source}{}{sourcepointer-1}{}}
\newlabel{algebra-lemma-weakly-ass-zero-divisors@cref}{{[originallemma][0][]05C3}{[1][0][]source}{}{}{}}
\makeatother
\crefname{originallemma}{original lemma}{original lemmas}
\Crefname{originallemma}{Original lemma}{Original lemmas}
\makeatletter
\newlabel{lemma-ass-weakly-ass}{{05AR}{source}{}{sourcepointer-2}{}}
\newlabel{lemma-ass-weakly-ass@cref}{{[originallemma][0][]05AR}{[1][0][]source}{}{}{}}
\makeatother
\crefname{originallemma}{original lemma}{original lemmas}
\Crefname{originallemma}{Original lemma}{Original lemmas}
\makeatletter
\newlabel{lemma-regular-section-structure-sheaf}{{01WZ}{source}{}{sourcepointer-3}{}}
\newlabel{lemma-regular-section-structure-sheaf@cref}{{[originallemma][0][]01WZ}{[1][0][]source}{}{}{}}
\makeatother
\crefname{originallemma}{original lemma}{original lemmas}
\Crefname{originallemma}{Original lemma}{Original lemmas}
\makeatletter
\newlabel{lemma-weakass-reduced}{{0EME}{source}{}{sourcepointer-4}{}}
\newlabel{lemma-weakass-reduced@cref}{{[originallemma][0][]0EME}{[1][0][]source}{}{}{}}
\makeatother
\title{Pullbacks of meromorphic functions are defined in seven situations ranging from flatness to weak association assumptions}
\author{The Stacks Project Authors}
\date{stacks-project:divisors.tex:local-snapshot-20261012; GNU FDL 1.2 or later}
\begin{document}
\maketitle
\noindent\textit{Editorial note.} All delivered mathematical text is the author's own, copied from the pinned Stacks Project chapter file pinned-original-source.tex; no statement or proof step was rewritten, reordered or removed.\par\medskip
\noindent\textit{Editorial note.} Following the multi-claim convention for this source, source\_url is the source-level URL https://stacks.math.columbia.edu/ and the source label lemma-pullback-meromorphic-sections-defined is carried in claim\_key and here rather than in a per-claim tag URL.\par\medskip
\noindent\textit{Editorial note.} The preamble is the chapter's shared Stacks Project preamble, retained as sources/stacks-divisors/retained-preamble.tex: a line-for-line excerpt of the shared Stacks preamble staged with this snapshot. No macro was invented or redefined, and this card adds no preamble text of its own.\par\medskip
\noindent\textit{Editorial note.} Statement span 5983-6004 and proof span 6006-6031 come from the halfA pool card operations/pool-cards/stacks-divisors-halfA.json and were re-read in the pinned chapter file before this card was written.\par\medskip
\noindent\textit{Editorial note.} The statement is a seven-item list; all seven cases are part of the counted statement span and the proof labels them by number exactly as the author does.\par\medskip
\noindent\textit{Editorial note.} 4 cross-references occur inside the counted spans; they are listed in external\_references as navigation pointers only.\par\medskip
\noindent\textit{Allowed standard prerequisites (editorial).}\begin{itemize}
\item The characterization of regular sections of the structure sheaf, which reduces the existence of pullbacks of meromorphic functions to the requirement that the ring map send nonzerodivisors to nonzerodivisors. Cited at proof lines 6008-6011. Turns the sheaf-theoretic statement into the algebraic statement about nonzerodivisors.
\item An element is a nonzerodivisor exactly when it lies outside every weakly associated prime. Cited at proof lines 6017-6020. Gives the main case, where the hypothesis on weakly associated points prevents the image of a nonzerodivisor from being a zerodivisor.
\item The comparison between associated and weakly associated primes, and the description of weakly associated points of a reduced scheme, both cited from the same chapter. Cited at proof lines 6021-6026. Derives the locally Noetherian cases and the reduced case from the main case.
\item Elements of a minimal prime of a ring are zerodivisors. Used at proof lines 6014-6016. Shows that a nonzerodivisor avoids every minimal prime, which is what the hypotheses on the images of the relevant points supply.
\end{itemize}\medskip
\ifdefined\hypertarget\hypertarget{sourcepointer-1}{}\fi
\noindent\textit{Source reference.} Original lemma 05C3: focused-seeds/AGP/sources/stacks-snapshot-20261011/tags.txt, source lines 6906--6906. Same-chapter reference used in the retained author text at source line(s) 2624, 6025, 6106, 6127. It is supplied only as a navigation pointer so that the retained text compiles; the counted item delivers the author's own statement and the author's complete proof, and no part of the referenced item is copied in or substituted for the counted argument.\par\medskip
\ifdefined\hypertarget\hypertarget{sourcepointer-2}{}\fi
\noindent\textit{Source reference.} Original lemma 05AR: focused-seeds/AGP/sources/stacks-snapshot-20261011/tags.txt, source lines 6859--6859. Same-chapter reference used in the retained author text at source line(s) 875, 965, 1222, 2625, 6026, 6345. It is supplied only as a navigation pointer so that the retained text compiles; the counted item delivers the author's own statement and the author's complete proof, and no part of the referenced item is copied in or substituted for the counted argument.\par\medskip
\ifdefined\hypertarget\hypertarget{sourcepointer-3}{}\fi
\noindent\textit{Source reference.} Original lemma 01WZ: focused-seeds/AGP/sources/stacks-snapshot-20261011/tags.txt, source lines 2500--2500. Same-chapter reference used in the retained author text at source line(s) 2934, 6011. It is supplied only as a navigation pointer so that the retained text compiles; the counted item delivers the author's own statement and the author's complete proof, and no part of the referenced item is copied in or substituted for the counted argument.\par\medskip
\ifdefined\hypertarget\hypertarget{sourcepointer-4}{}\fi
\noindent\textit{Source reference.} Original lemma 0EME: focused-seeds/AGP/sources/stacks-snapshot-20261011/tags.txt, source lines 18304--18304. Same-chapter reference used in the retained author text at source line(s) 2628, 6029, 6463. It is supplied only as a navigation pointer so that the retained text compiles; the counted item delivers the author's own statement and the author's complete proof, and no part of the referenced item is copied in or substituted for the counted argument.\par\medskip
\begin{lemma}
\label{lemma-pullback-meromorphic-sections-defined}
Let $f : X \to Y$ be a morphism of schemes.
In each of the following cases pullbacks of meromorphic
functions are defined.
\begin{enumerate}
\item every weakly associated point of $X$ maps to
a generic point of an irreducible component of $Y$,
\item $X$, $Y$ are integral and $f$ is dominant,
\item $X$ is integral and the generic point of $X$ maps
to a generic point of an irreducible component of $Y$,
\item $X$ is reduced and every generic point of every irreducible
component of $X$ maps to the generic point of an irreducible component
of $Y$,
\item $X$ is locally Noetherian, and any associated point of
$X$ maps to a generic point of an irreducible component of $Y$,
\item $X$ is locally Noetherian, has no embedded points and
any generic point of an irreducible component of
$X$ maps to the generic point of an irreducible component of $Y$, and
\item $f$ is flat.
\end{enumerate}
\end{lemma}

\begin{proof}
The question is local on $X$ and $Y$. Hence we reduce to the case where
$X = \Spec(A)$, $Y = \Spec(R)$ and $f$ is given by a ring map
$\varphi : R \to A$.
By the characterization of regular sections of the structure sheaf
in Lemma \ref{lemma-regular-section-structure-sheaf} we have to
show that $R \to A$ maps nonzerodivisors to nonzerodivisors.
Let $t \in R$ be a nonzerodivisor.

\medskip\noindent
If $R \to A$ is flat, then $t : R \to R$ being injective
shows that $t : A \to A$ is injective. This proves (7).

\medskip\noindent
In the other cases we note that $t$ is not contained in any of
the minimal primes of $R$ (because every element of a minimal
prime in a ring is a zerodivisor).
Hence in case (1) we see that $\varphi(t)$ is not contained
in any weakly associated prime of $A$. Thus this case follows from
Algebra, Lemma \ref{algebra-lemma-weakly-ass-zero-divisors}.
Case (5) is a special case of (1) by Lemma \ref{lemma-ass-weakly-ass}.
Case (6) follows from (5) and the definitions.
Case (4) is a special case of (1) by
Lemma \ref{lemma-weakass-reduced}.
Cases (2) and (3) are special cases of (4).
\end{proof}

\end{document}
